\chapter*{U}
\label{chap:u}

\term{Unit}
A unit is either the multiplicative identity in a ring or an element that has a multiplicative inverse, depending on context. \par\noindent\textbf{Applications:} units describe reversible scaling and invertibility in algebra.

\term{Unit (Ring Theory)}
A unit of a ring $R$ is an element $u$ with an inverse $v$ satisfying $uv=vu=1$. The units form the group $R^\times$; in $\mathbb Z$ they are $1$ and $-1$. \par\noindent\textbf{Applications:} units identify invertible coefficients and symmetries in algebra and number theory.

\term{Unit Circle}
The unit circle is the circle $x^2+y^2=1$ centered at the origin. Its points are $(\cos\theta,\sin\theta)$. \par\noindent\textbf{Applications:} it defines trigonometric functions and models rotations, periodic motion, and complex numbers.

\term{Unit Group}
The unit group $R^\times$ of a ring consists of all elements with multiplicative inverses. In a field it is the group of all nonzero elements. \par\noindent\textbf{Applications:} unit groups encode invertibility in algebraic number theory and modular arithmetic.

\term{Unit Vector}
A unit vector has norm one. A nonzero vector $v$ can be normalized as $\hat v=v/\|v\|$. \par\noindent\textbf{Applications:} unit vectors represent directions, orientations, and normalized features.

\term{Unitary Matrix}
A unitary matrix satisfies $U^*U=UU^*=I$. It preserves inner products, lengths, and angles in complex space. \par\noindent\textbf{Applications:} unitary matrices describe quantum gates, rotations, Fourier transforms, and stable changes of basis.

\term{Unitary Operator}
A unitary operator is a surjective linear operator on a Hilbert space that preserves inner products. It is the infinite-dimensional analogue of a unitary matrix. \par\noindent\textbf{Applications:} unitary operators model quantum evolution and energy-preserving transformations.

\term{Universal}
Universal means applying to every case in a specified domain. The symbol $\forall$ expresses “for all.” \par\noindent\textbf{Applications:} universal statements formulate definitions, theorems, specifications, and logical rules.

\term{Universal Set}
A universal set is the chosen set containing every object under discussion; all sets in that context are its subsets. \par\noindent\textbf{Applications:} it fixes the complement operation and organizes elementary set calculations.

\term{Universal Mapping Property}
A universal mapping property characterizes an object by a unique map or factorization property rather than by a formula. Tensor products and free groups are examples. \par\noindent\textbf{Applications:} it gives canonical constructions throughout algebra, topology, and category theory.

\term{Uncountable Set}
An uncountable set is an infinite set that cannot be put in bijection with $\mathbb N$. The real numbers are uncountable. \par\noindent\textbf{Applications:} uncountability distinguishes sizes of infinity and limits enumeration and computation.

\term{Undecidable}
An undecidable problem has no algorithm that always halts and returns the correct answer for every input. \par\noindent\textbf{Applications:} undecidability identifies fundamental limits of software verification and computation.

\term{Underdetermined}
An underdetermined system has fewer independent equations than unknowns and usually has many solutions. \par\noindent\textbf{Applications:} underdetermined models arise in inverse problems, data fitting, and parameter estimation.

\term{Uniform}
Uniform means that one bound or choice works throughout a whole domain, rather than changing from point to point. \par\noindent\textbf{Applications:} uniform estimates control approximation, convergence, and numerical error.

\term{Uniform Continuity}
Uniform continuity means that for every $\epsilon>0$ one $\delta>0$ works for all points in the domain. Unlike ordinary continuity, $\delta$ does not depend on the point. \par\noindent\textbf{Applications:} it preserves controlled approximation and ensures uniform limits behave well.

\term{Uniform Convergence}
Uniform convergence means that the largest error between functions and their limit tends to zero. \par\noindent\textbf{Applications:} it allows limits to preserve continuity and supports reliable approximation by function sequences.

\term{Uniform Boundedness Principle}
See \term{Banach--Steinhaus Theorem}. \par\noindent\textbf{Applications:} it controls families of linear operators using pointwise bounds.

\term{Uniform Distribution}
A uniform distribution assigns equal probability to equal-sized regions of its support. On $[a,b]$ its density is $1/(b-a)$. \par\noindent\textbf{Applications:} it models unbiased random choices and supplies a basic simulation distribution.

\term{Uniform Space}
A uniform space is a set with a structure describing when pairs of points are uniformly close, even when no distance function is given. \par\noindent\textbf{Applications:} uniform spaces define completeness and uniform continuity in general topology.

\term{Union}
The union of sets contains every element belonging to at least one set: $A\cup B=\{x:x\in A\text{ or }x\in B\}$. \par\noindent\textbf{Applications:} unions combine alternatives and describe events, regions, and solution sets.

\term{Unimodular Matrix}
A unimodular matrix has integer entries and determinant $\pm1$, so its inverse also has integer entries. \par\noindent\textbf{Applications:} unimodular matrices preserve integer lattices and support integer programming and lattice algorithms.

\term{Uniqueness Theorem}
A uniqueness theorem gives conditions under which a problem has at most one solution. \par\noindent\textbf{Applications:} uniqueness establishes well-defined models and is central to differential equations, optimization, and inverse problems.

\term{Unipotent Element}
A unipotent element satisfies that $u-1$ is nilpotent. A unipotent matrix has only eigenvalue $1$ and is similar to an upper triangular matrix with ones on the diagonal. \par\noindent\textbf{Applications:} unipotent elements appear in algebraic groups, Lie theory, and matrix decompositions.

\term{Upper Half-Plane}
The upper half-plane is $\{z\in\mathbb C:\operatorname{Im}(z)>0\}$. It is a standard domain for complex and hyperbolic geometry. \par\noindent\textbf{Applications:} it supports modular forms, conformal maps, and the Poincaré model.

\term{Upper Bound}
An upper bound $M$ of a set $S$ satisfies $s\leq M$ for every $s\in S$. The least upper bound is its supremum. \par\noindent\textbf{Applications:} upper bounds control errors, sizes, optimization objectives, and convergence.

\term{Upper Semi-continuous Function}
An upper semi-continuous function cannot jump upward in the limit: nearby values have limit superior at most the value at the point. \par\noindent\textbf{Applications:} upper semi-continuity helps prove existence of maxima and stable optimization solutions.

\term{Upper Triangular}
An upper-triangular matrix has zeros below its main diagonal. Its eigenvalues are the diagonal entries. \par\noindent\textbf{Applications:} triangular form simplifies eigenvalue calculation, linear systems, and matrix exponentials.

\term{Urysohn Lemma}
The Urysohn lemma states that disjoint closed sets in a normal space can be separated by a continuous function taking value $0$ on one and $1$ on the other. \par\noindent\textbf{Applications:} it constructs partitions of unity and supports metrization and extension theorems.

\term{Urysohn Metrization Theorem}
The Urysohn metrization theorem states that every second-countable regular space admits a compatible metric. \par\noindent\textbf{Applications:} it converts abstract topological questions into metric and analytical ones.

\term{Use}
Use is the act of applying a function, rule, or definition to given arguments. \par\noindent\textbf{Applications:} precise use of rules supports symbolic calculation, programming, and formal proofs.

\term{U-statistic}
A U-statistic is an unbiased estimator formed by averaging a kernel over subsets of a sample. \par\noindent\textbf{Applications:} U-statistics estimate population quantities such as moments, correlations, and interaction effects.

\term{Unitary Group}
The unitary group $U(n)$ consists of all complex $n\times n$ unitary matrices. It is a compact Lie group. \par\noindent\textbf{Applications:} it describes quantum symmetries, rotations of complex state spaces, and harmonic analysis.

\term{Unique}
Unique means existing in exactly one form or instance. \par\noindent\textbf{Applications:} uniqueness identifies a determined solution, canonical object, or unambiguous representation.

\term{Unique Factorization Domain (UFD)}
A unique factorization domain is an integral domain in which every nonzero nonunit factors uniquely into irreducible, equivalently prime, elements up to order and units. \par\noindent\textbf{Applications:} UFDs support polynomial factorization and arithmetic arguments.

\term{Unital Ring}
A unital ring has a multiplicative identity $1$ satisfying $1a=a1=a$ for every element $a$. \par\noindent\textbf{Applications:} unital rings provide the standard setting for modules, algebras, matrices, and algebraic geometry.

\term{Unitary Representation}
A unitary representation maps a group to unitary operators on a Hilbert space. \par\noindent\textbf{Applications:} it represents symmetries in quantum mechanics, harmonic analysis, and signal theory.

\term{Unimodular Group}
A locally compact group is unimodular when its left and right Haar measures agree. \par\noindent\textbf{Applications:} unimodularity simplifies integration, harmonic analysis, and representation theory on groups.

\term{Universal Covering Space}
A universal covering space is a simply connected space mapping locally like a given space through a covering map. \par\noindent\textbf{Applications:} it reveals fundamental groups and supports classification of covering spaces and manifolds.

\term{Ultrafilter}
An ultrafilter is a maximal collection of subsets closed under finite intersections and supersets, containing exactly one of each set and its complement. \par\noindent\textbf{Applications:} ultrafilters support compactness proofs, limits, and nonstandard analysis.

\term{Ultraproduct}
An ultraproduct combines a family of structures and identifies sequences that agree on a chosen ultrafilter-large set. \par\noindent\textbf{Applications:} ultraproducts transfer first-order properties and support model theory and functional analysis.

\term{Ultrametric Space}
An ultrametric space satisfies $d(x,z)\leq\max\{d(x,y),d(y,z)\}$, a stronger triangle inequality than for ordinary metrics. \par\noindent\textbf{Applications:} ultrametrics model $p$-adic numbers, hierarchical clustering, and phylogenetic trees.

\term{Unary Operation}
A unary operation takes one input and returns one output, such as negation or complementation. \par\noindent\textbf{Applications:} unary operations define algebraic structures, logic, syntax, and programming expressions.

\term{Unconditional Convergence}
A series converges unconditionally when every rearrangement converges to the same sum. In many important spaces this is equivalent to absolute convergence. \par\noindent\textbf{Applications:} unconditional convergence supports rearranging infinite expansions and stable functional analysis.

\term{Unbounded Operator}
An unbounded operator is a linear operator that has no finite global norm bound on its domain. \par\noindent\textbf{Applications:} unbounded operators model derivatives, momentum, and energy in PDEs and quantum mechanics.

\term{Unit Disc}
The unit disc is $\{z\in\mathbb C:|z|\leq1\}$; its interior is $|z|<1$. \par\noindent\textbf{Applications:} it is a standard domain for complex analysis, conformal maps, and power series.

\term{Unitary Transformation}
A unitary transformation preserves inner products and therefore lengths and angles in a complex vector space. \par\noindent\textbf{Applications:} unitary transformations represent quantum evolution, rotations, and Fourier changes of basis.

\term{Roots of Unity}
The $n$th roots of unity are the complex numbers satisfying $z^n=1$, namely $e^{2\pi ik/n}$ for $k=0,\ldots,n-1$. \par\noindent\textbf{Applications:} they support Fourier analysis, cyclic symmetry, polynomial factorization, and coding.

\term{Univalent Function}
A univalent function is a holomorphic function that is one-to-one on its domain. \par\noindent\textbf{Applications:} univalent functions control conformal mappings and geometric function theory.
